% ============================================
% Johns Hopkins Letter Template
% Uses the built-in 'letter' class
% ============================================

\documentclass[11pt,letterpaper]{letter}

\usepackage{graphicx}
\usepackage{eso-pic}
\usepackage{geometry}
\usepackage{microtype}
\usepackage[utf8]{inputenc}
\usepackage[hidelinks]{hyperref}

% ----- Page Setup -----
\geometry{left=25mm,right=25mm,top=25mm,bottom=25mm}

% ----- Logo Placement (foreground, first page only) -----
\newcommand{\PutJHULogoFG}[1]{%
  \AddToShipoutPictureFG*{%
    \AtPageUpperLeft{%
      \hspace{10mm}%
      \raisebox{-38mm}{%
        \includegraphics[width=6cm]{#1}%
      }%
    }%
  }%
}

% ----- Sender Information -----
\signature{Decory Edwards}
\address{
Department of Economics \\
Johns Hopkins University \\
Baltimore, MD 21218 \\
\texttt{dedwar65@jh.edu}
}

\begin{document}
\PutJHULogoFG{Images/jhulogo.png}

\begin{letter}{
Search Committee \\
Department of Economics and Accounting \\
College of the Holy Cross \\
Worcester, MA
}

\opening{Dear Members of the Search Committee,}

I am writing to apply for the visiting full-time faculty position in Economics at the College of the Holy Cross for the 2026–27 academic year. I am a Ph.D.\ candidate in the Department of Economics at Johns Hopkins University (advisors: Christopher D.\ Carroll, Nicholas W.\ Papageorge, and Stelios Fourakis), and I expect to receive my degree in May 2026. My primary specializations are macroeconomic modeling, wealth and income dynamics, and computational economics.

My research agenda focuses on the ability of macroeconomic models to generate empirically realistic distributions of wealth and income over the life cycle. A key driver of that work is modeling heterogeneity in the rate of return on assets, and understanding how return dispersion contributes to portfolio accumulation across households.

In my job market paper, I calibrate a life cycle heterogeneous agent model to match wealth moments from the Survey of Consumer Finances by allowing households to differ in their portfolio returns. The model helps explain observed wealth concentration and return dispersion. In a related research project using the Health and Retirement Study, I examine how trust influences both income growth and asset returns. I find a hump shaped relationship for trust and returns (consistent with the literature) and characterize the distribution of the persistent component of returns to net worth. These experiences have equipped me to frame research strategies and design modeling approaches, execute econometric and simulation analysis with large micro data sets, and translate results into clear and rigorous written deliverables for both academic and practitioner audiences.

I am especially drawn to Holy Cross because of its strong commitment to undergraduate teaching within a Jesuit liberal arts tradition. Teaching Principles of Economics provides an opportunity to introduce students to economic reasoning as a tool for understanding both markets and public policy, while emphasizing ethical reflection and the social dimensions of economic life. I have experience serving as a teaching assistant in quantitative and applied courses and have led review sessions and discussions that focus on clarity, accessibility, and active engagement. I would be enthusiastic about teaching Principles and Statistics, and potentially offering an elective related to macroeconomics, inequality, or public finance. The opportunity to teach three courses per semester in a community that values close faculty student interaction aligns closely with my professional goals.

Holy Cross’s mission to educate students in the service of others resonates strongly with my interest in the distributional consequences of economic policy. My research on taxation and heterogeneity connects directly to questions about economic opportunity, institutional design, and social responsibility. I am committed to inclusive pedagogy that supports students from diverse backgrounds and fosters respectful dialogue across differences. In the classroom, I aim to create an environment where students feel comfortable engaging with challenging material and developing analytical tools that serve them beyond college.

I am enthusiastic about relocating to Worcester and teaching in person as part of a residential academic community. A visiting appointment at Holy Cross would provide an ideal setting to further develop both my teaching and research profile, while contributing meaningfully to the intellectual life of the department.

I believe I would be a strong fit for this position because:
\begin{enumerate}
    \item My research background equips me to bring contemporary macroeconomic and public finance insights into the undergraduate classroom.
    \item My teaching experience emphasizes clarity, engagement, and inclusive pedagogical practices.
    \item I value the liberal arts tradition and the integration of analytical rigor with reflection on ethical and social questions.
\end{enumerate}

Thank you very much for your time and consideration. I would welcome the opportunity to contribute to the Department of Economics and Accounting at the College of the Holy Cross.

\closing{Sincerely,}

\end{letter}
\end{document}
